\documentclass[11pt,a4paper]{report}
\usepackage{coursmodern}

% 1. Paramétrage de la page de garde
\CourseTitle{Espaces Vectoriels \& Dimension Finie}
\CourseSubtitle{EMINES --- CPI2A (Class2030)}
\ChapterNumber{01}
\CourseAuthor{Samy Youssoufine}
\LastUpdate{\today}
\PrimarySite{maths2.txt.ma}{https://maths2.txt.ma}
\MirrorSite{samy-y.github.io/maths-cpi2a}{https://samy-y.github.io/maths-cpi2a}

% Activer ou commenter pour masquer la boîte de chantier
\wiptrue

\begin{document}

\MakeSwissCover

\tableofcontents

\chapter{Groupes symétriques}

\section{Permutations}

\begin{definition}[Permutation et ensemble de permutations]
    Soit $n \in \N^*$. Une \textbf{permutation} de $\{1, 2, \dots, n\}$ est une bijection de cet ensemble dans lui-même.

    On écrit :
    $$\sigma = \left(\begin{matrix} 1 & 2 & \dots & n \\ \sigma(1) & \sigma(2) & \dots & \sigma(n) \end{matrix}\right)$$
    pour représenter la permutation $\sigma$.

    On note $S_n$ l'ensemble de toutes les permutations de $\{1, 2, \dots, n\}$.
\end{definition}

\begin{example}[Permutation]
    Par exemple, la permutation $\sigma$ suivante :
    \[
        \begin{array}{c|ccccc}
            i & 1 & 2 & 3 & 4 & 5 \\
            \hline
            \sigma(i) & 3 & 1 & 5 & 4 & 2
        \end{array}
    \]
    $\sigma$ est bel et bien une bijection de $\{1, 2, 3, 4, 5\}$ dans lui-même.
\end{example}

\begin{example}[Ensemble de permutations]
    $S_2$ contient les permutations suivantes :
    \[
        S_2 = \left\{\left(\begin{matrix} 1 & 2 \\ 1 & 2 \end{matrix}\right), \left(\begin{matrix} 1 & 2 \\ 2 & 1 \end{matrix}\right)\right\}
    \]
    $S_3$ contient les permutations suivantes :
    \[
        S_3 = \left\{\left(\begin{matrix} 1 & 2 & 3 \\ 1 & 2 & 3 \end{matrix}\right), \left(\begin{matrix} 1 & 2 & 3 \\ 1 & 3 & 2 \end{matrix}\right), \left(\begin{matrix} 1 & 2 & 3 \\ 2 & 1 & 3 \end{matrix}\right), \left(\begin{matrix} 1 & 2 & 3 \\ 2 & 3 & 1 \end{matrix}\right), \left(\begin{matrix} 1 & 2 & 3 \\ 3 & 1 & 2 \end{matrix}\right), \left(\begin{matrix} 1 & 2 & 3 \\ 3 & 2 & 1 \end{matrix}\right)\right\}
    \]
\end{example}

\begin{definition}[Composée de permutations]
    Soient $\sigma_1$ et $\sigma_2$ deux permutations de $S_n$. La \textbf{composée} de $\sigma_1$ et $\sigma_2$, notée $\sigma_1 \circ \sigma_2$, est définie par :
    $$\forall i \in [\![1,n]\!], (\sigma_1 \circ \sigma_2)(i) = \sigma_1(\sigma_2(i))$$
\end{definition}

\begin{example}[Composée de deux permutations]
    Soient $\sigma_1$ et $\sigma_2$ les permutations définies par :
    $$
    \begin{cases}
        \sigma_1 = \left(\begin{matrix} 1 & 2 & 3 & 4 & 5 \\ 5 & 2 & 3 & 1 & 4 \end{matrix}\right) \\
        \sigma_2 = \left(\begin{matrix} 1 & 2 & 3 & 4 & 5 \\ 4 & 3 & 5 & 1 & 2 \end{matrix}\right)
    \end{cases}
    $$
    La composée $\sigma_1 \circ \sigma_2$ est définie par :
    $$
    (\sigma_1 \circ \sigma_2)(i) = \sigma_1(\sigma_2(i))
    $$
    Essayons pour toutes les valeurs de $i$ :
    \[
    \begin{array}{c|ccccc}
        i & 1 & 2 & 3 & 4 & 5 \\
        \hline
        \sigma_2(i) & 4 & 3 & 5 & 1 & 2 \\
        \sigma_1(\sigma_2(i)) & \sigma_1(4) & \sigma_1(3) & \sigma_1(5) & \sigma_1(1) & \sigma_1(2) \\
        & 1 & 3 & 4 & 5 & 2
    \end{array}
    \]
    Ainsi, la composée $\sigma_1 \circ \sigma_2$ est donnée par :
    $$
    \sigma_1 \circ \sigma_2 = \left(\begin{matrix} 1 & 2 & 3 & 4 & 5 \\ 1 & 3 & 4 & 5 & 2 \end{matrix}\right)
    $$
    Notons aussi que la composée $\sigma_2 \circ \sigma_1$ est donnée par :
    $$
    \sigma_2 \circ \sigma_1 = \left(\begin{matrix} 1 & 2 & 3 & 4 & 5 \\ 4 & 3 & 5 & 1 & 2 \end{matrix}\right)
    $$
    Les deux composées sont différentes, ce qui montre que la composition de permutations n'est pas commutative.
\end{example}

\begin{definition}[Inverse d'une permutation]
    Soit $\sigma$ une permutation de $S_n$. L'\textbf{inverse} de $\sigma$, notée $\sigma^{-1}$, est la permutation définie par :
    $$\forall i \in [\![1,n]\!], \sigma^{-1}(\sigma(i)) = i$$
\end{definition}

\begin{example}[Inverse d'une permutation]
    Réutilisons la permutation $\sigma_1$ définie par :
    $$
\sigma_1 = \left(\begin{matrix} 1 & 2 & 3 & 4 & 5 \\ 5 & 2 & 3 & 1 & 4 \end{matrix}\right)
    $$
    L'inverse de $\sigma_1$, noté $\sigma_1^{-1}$, est défini par :
    $$
\sigma_1^{-1}(\sigma_1(i)) = i
    $$
    Essayons pour toutes les valeurs de $i$ :
    \[
\begin{array}{c|ccccc}
    i & 1 & 2 & 3 & 4 & 5 \\
    \sigma_1^{-1}(i) & 4 & 2 & 3 & 5 & 1
\end{array}
    \]
    Ainsi, l'inverse de $\sigma_1$ est donné par :
    $$
\sigma_1^{-1} = \left(\begin{matrix} 1 & 2 & 3 & 4 & 5 \\ 4 & 2 & 3 & 5 & 1 \end{matrix}\right)
    $$
\end{example}

\begin{property}
    Il y a $n!$ permutations de $\{1, 2, \dots, n\}$. Le cardinal de $S_n$ est donc $n!$.

    $$\card{S_n} = n!$$
\end{property}

\begin{proof}
    Chaque élément de $\{1, 2, \dots, n\}$ doit être envoyé vers un élément distinct de l'ensemble d'arrivée. Pour le premier élément, il y a $n$ choix possibles, pour le deuxième élément, il y a $n-1$ choix possibles (car on ne peut pas réutiliser l'élément déjà choisi pour le premier), pour le troisième élément, il y a $n-2$ choix possibles, et ainsi de suite jusqu'au dernier élément qui n'a plus qu'un seul choix possible.
\end{proof}

\begin{property}
    \begin{enumerate}
        \item \textbf{Associativité} : Pour toutes permutations $\sigma_1$, $\sigma_2$ et $\sigma_3$ de $S_n$, on a :
        $$(\sigma_1 \circ \sigma_2) \circ \sigma_3 = \sigma_1 \circ (\sigma_2 \circ \sigma_3)$$
        On hérite ces propriétés de l'associativité de la composition de fonctions.
        \item \textbf{Existence d'un élément neutre} : Il existe une permutation $\Id$ dans $S_n$ telle que pour toute permutation $\sigma$ de $S_n$, on a :
        $$\Id \circ \sigma = \sigma \circ \Id = \sigma$$
        L'élément neutre est la permutation identité, notée $\Id$, définie par :
        $$\Id = \left(\begin{matrix} 1 & 2 & \dots & n \\ 1 & 2 & \dots & n \end{matrix}\right)$$
        \item \textbf{Existence d'inverses} : Pour toute permutation $\sigma$ de $S_n$, il existe une permutation $\sigma^{-1}$ de $S_n$ telle que :
        $$\sigma \circ \sigma^{-1} = \sigma^{-1} \circ \sigma = \Id$$
    \end{enumerate}
\end{property}

\chapter{Bases et Dimension Finie}

\section{Théorèmes fondamentaux}

\begin{definition}[Famille génératrice finie]
Un espace vectoriel $E$ est dit de dimension finie s'il admet une famille génératrice finie.
\end{definition}

\todo{Recopier la démonstration de l'extraction d'une base libre.}

\begin{theorem}[Théorème de la base incomplète]
Soit $E$ un $\mathbb{K}$-espace vectoriel de dimension finie $n$. Toute famille libre de vecteurs de $E$ peut être complétée en une base de $E$.
\end{theorem}

\begin{property}
Si $F$ et $G$ sont deux sous-espaces de $E$ avec $F \subset G$ et $\dim(F) = \dim(G) < +\infty$, alors $F = G$.
\end{property}

\begin{consequence}
Dans un espace de dimension $n$, toute famille libre de $n$ vecteurs est une base.
\end{consequence}

\begin{corollary}[Isomorphismes]
Un endomorphisme $u \in \mathcal{L}(E)$ est injectif si et seulement s'il est surjectif.
\end{corollary}

\begin{application}
Vérification de l'inversibilité d'une matrice carrée par inverse unilatéral.
\end{application}

\begin{remark}
Par convention, l'espace nul $\{0_E\}$ a pour dimension 0.
\end{remark}

\begin{example}
L'espace $\mathbb{R}_n[X]$ des polynômes de degré au plus $n$ est de dimension $n+1$.
\end{example}

\begin{attention}[Dimension infinie]
\label{attention:dim_inf}
La relation $\dim(F + G) = \dim F + \dim G - \dim(F \cap G)$ nécessite la dimension finie.
\end{attention}

% Exercice avec niveau de difficulté 2 étoiles et titre externe
\begin{exercise}[2][Noyaux itérés de nilpotence]
\label{exo:noyaux_iteres}
Soit $u \in \mathcal{L}(E)$ avec $\dim(E) = n$. Montrer que la suite $(\ker(u^k))_{k \in \mathbb{N}}$ stationne au plus tard au rang $n$.
\end{exercise}

% Correction avec référence externe
\begin{correction}[Exercice \ref{exo:noyaux_iteres}]
On étudie la stricte inclusion des noyaux successifs avant le rang de stationnement.
\end{correction}

\begin{methodbox}[Pivot de Gauss]
Échelonnement en colonnes pour extraire une sous-famille libre maximale.
\end{methodbox}

% 13e environnement libre : choix direct de l'icône FontAwesome, titre, couleur, argument
\begin{customenv}{\faMicrochip}{Algorithmique}[cMeth][Décomposition LU]
En calcul numérique matriciel, l'algorithme correspond directement à l'élimination de Gauss sans pivotage partiel.
\end{customenv}

% -----------------------------------------------------------------------
% FIGURE DE TEST TIKZ 3D
% -----------------------------------------------------------------------
\section{Figure de test TikZ 3D}

\begin{example}[Repère cartésien 3D et cube unité]
La figure ci-dessous est un test de rendu TikZ 3D utilisant le paquet \texttt{tikz-3dplot}.
Elle affiche un repère cartésien orienté, un cube unité aux faces semi-transparentes, et
un vecteur diagonal $\vec{v} = \frac{1}{\sqrt{3}}(1,1,1)$.
\end{example}

\begin{figure}[h]
  \centering
  % Angles de vue : inclinaison 68°, rotation 115°
  \tdplotsetmaincoords{68}{115}
  \begin{tikzpicture}[tdplot_main_coords, scale=2.8,
      axis/.style   = {thick, ->, >=Stealth},
      edge3d/.style = {thick},
      face/.style   = {opacity=0.12},
    ]

    %--- Faces du cube (dessinées avant les arêtes pour la transparence) ---
    % Face inférieure  (z=0)
    \fill[face, fill=cDef]   (0,0,0)--(1,0,0)--(1,1,0)--(0,1,0)--cycle;
    % Face supérieure  (z=1)
    \fill[face, fill=cThm]   (0,0,1)--(1,0,1)--(1,1,1)--(0,1,1)--cycle;
    % Face avant       (y=0)
    \fill[face, fill=cProp]  (0,0,0)--(1,0,0)--(1,0,1)--(0,0,1)--cycle;
    % Face arrière     (y=1)
    \fill[face, fill=cProp]  (0,1,0)--(1,1,0)--(1,1,1)--(0,1,1)--cycle;
    % Face gauche      (x=0)
    \fill[face, fill=cCor]   (0,0,0)--(0,1,0)--(0,1,1)--(0,0,1)--cycle;
    % Face droite      (x=1)
    \fill[face, fill=cCor]   (1,0,0)--(1,1,0)--(1,1,1)--(1,0,1)--cycle;

    %--- Arêtes du cube ---
    % Base
    \draw[edge3d, cDef!70!black]
      (0,0,0)--(1,0,0)--(1,1,0)--(0,1,0)--cycle;
    % Sommet
    \draw[edge3d, cThm!70!black]
      (0,0,1)--(1,0,1)--(1,1,1)--(0,1,1)--cycle;
    % Montants verticaux
    \foreach \x/\y in {0/0, 1/0, 1/1, 0/1}{
      \draw[edge3d, swissgray] (\x,\y,0)--(\x,\y,1);
    }

    %--- Axes cartésiens ---
    \draw[axis, cDef]   (0,0,0) -- (1.6,0,0) node[anchor=north east,  font=\sffamily\bfseries\small]{$x$};
    \draw[axis, cProp]  (0,0,0) -- (0,1.6,0) node[anchor=north west,  font=\sffamily\bfseries\small]{$y$};
    \draw[axis, cThm]   (0,0,0) -- (0,0,1.6) node[anchor=south,       font=\sffamily\bfseries\small]{$z$};

    % Origine
    \fill[swissblack] (0,0,0) circle (0.018) node[below left, font=\sffamily\tiny]{$O$};

    %--- Vecteur diagonal ---
    \draw[->, >=Stealth, cAtt, ultra thick, dashed]
      (0,0,0) -- (0.75,0.75,0.75)
      node[right, font=\sffamily\small, text=cAtt]{$\vec{v}$};

    % Projections pointillées du vecteur
    \draw[cAtt!40, thin, dotted] (0.75,0.75,0.75)--(0.75,0.75,0);
    \draw[cAtt!40, thin, dotted] (0.75,0.75,0.75)--(0.75,0,0.75);
    \draw[cAtt!40, thin, dotted] (0.75,0.75,0.75)--(0,0.75,0.75);

  \end{tikzpicture}
  \caption{%
    Test TikZ 3D --- Repère cartésien $(O, \vec{x}, \vec{y}, \vec{z})$,
    cube unité aux faces colorées semi-transparentes,
    et vecteur diagonal $\vec{v}$.%
  }
  \label{fig:test3d}
\end{figure}

\end{document}